\chapter{Terminology reference}
\label{chap:terminos}
\begin{description}
 \item[Backend] Implementation selected to validate or solve a model family.
 The validation option does not compute a solution.
 \item[Balance] Equality accounting for inflows, outflows, and changes in state,
 such as ending inventory equaling initial inventory plus production minus demand.
 \item[Capacity] Resource limit, with a defined unit and time horizon.
 \item[Constraint] Condition limiting the admissible decisions.
 \item[Decision variable] Controllable quantity chosen by the decision maker in the model.
 \item[Diagnostic] Validation message with severity, code, and, when available,
 the path of the affected field.
 \item[Exactness] Characterization of the procedure and its guarantees for the
 supported model, distinct from the number of decimal places displayed.
 \item[Expected value] Probability-weighted average; it does not guarantee that
 payoff in any particular realization.
 \item[Feasibility] Satisfaction of the model's constraints, bounds, and
 conditions; it does not imply optimality.
 \item[Feasible region] Set of solutions satisfying the constraints.
 \item[Global optimum] Feasible solution with the best objective throughout the
 feasible region; a local optimum is compared only with nearby solutions.
 \item[Heuristic] Procedure seeking a good solution without necessarily
 guaranteeing a global optimum.
 \item[Likelihood] Probability of observing evidence conditional on a state;
 it is not the posterior probability of that state.
 \item[Linear relaxation] Model replacing integrality requirements with
 continuous variables; its solution may be inadmissible for the integer problem.
 \item[Makespan] Completion time of the last job in a schedule.
 \item[Normalized JSON] Structured representation of a model or result that
 allows data to be exchanged, inspected, and reproduced across QSB clients.
 \item[Objective function] Expression quantifying the criterion to maximize or minimize.
 \item[Optimal solution] Feasible solution achieving the best objective of the
 problem considered, within the method's guarantees and tolerances.
 \item[Parameter] Data item treated as known by the model during a run.
 \item[Policy] Rule or set of decisions by stage and state.
 \item[Posterior] Probability of a state after observing evidence, calculated
 from priors and likelihoods.
 \item[Seed] Value initializing the pseudorandom generator; it enables a
 simulation to be repeated with the same model, configuration, and implementation.
 \item[Sensitivity] Analysis of how the solution changes when data or assumptions change.
 \item[Shadow price] Marginal change in the objective when a resource varies
 within the range and under the conditions where the dual interpretation is valid.
 \item[Slack] Remaining capacity or margin of a constraint for a solution.
 \item[State] Summary of sufficient information for future decisions in a
 dynamic model; it may also refer to a random condition of the environment.
 \item[Tolerance] Numerical threshold used to compare quantities or accept
 approximations; it must be interpreted alongside the scale of the problem.
 \item[Trace] Record of stages, states, and decisions helping to examine a
 computed policy; it is not necessarily the algorithm's entire internal table.
 \item[Utilization] Fraction of service capacity occupied on average.
 \item[Value of information] Expected improvement from allowing decisions based
 on evidence; EVSI refers to sample information and EVPI to perfect information.
\end{description}
